\begin{claim}{D.2.60}{Proved}{The auxiliary James--Hopf class survives all twelve-cells}
Put $B=K(\mathbb F_2,2)$, $X_0=\Sigma B$, $Y=\Sigma X_0$, and
\[
 Z=\Sigma(X_0\wedge X_0)=\Sigma^3(B\wedge B).
 \tag{D.2.60a}
\]
Work at the prime two and use the actual CW model specified below.
If $i_{7,12}:S^7\to Z^{(12)}$ is its bottom map, then
\[
 0\ne i_{7,12}\sigma'\in\pi_{14}Z^{(12)},\qquad
                         2i_{7,12}\sigma'=0.
 \tag{D.2.60b}
\]
Its image in $\pi_{14}Z$ is the actual second James--Hopf image
of $\alpha\sigma'$, where $\alpha=u_4\nu_4\in\pi_7Y$.
Thus none of the ten-, eleven-, or twelve-dimensional attaching
maps kills the indicated auxiliary class.  This assertion does
not determine its image after the cells in dimensions thirteen
through fifteen, or the value of $\alpha\sigma'$ or $h_Ml_4$ in
the full target.
\end{claim}
\begin{proof}
All free groups and cellular degrees below are two-local.
Write $M(2,j)=M(\mathbb Z/2,j)$ for a Moore space with bottom
cell in degree $j$, and put
\[
 C_1^j=\operatorname{Cof}
 \bigl((2,\eta_{j-2}):S^{j-2}\vee S^{j-1}\longrightarrow S^{j-2}\bigr).
 \tag{D.2.60c}
\]
We retain actual attaching maps when their composites leave
the stable range.

\textbf{The bottom class and the borderline nine-stage.}
Naturality of the second James--Hopf map with respect to
$\Sigma u_3:S^4\to Y$, together with the sphere calculation
$H(\nu_4\sigma')=\sigma'$ in
\cite[Proposition~13, p.~9]{FGHStricklandToda}, gives
\[
 H_2(\alpha\sigma')=i_7\sigma',\qquad
                  i_7=\Sigma(u_3\wedge u_3):S^7\to Z.
 \tag{D.2.60d}
\]
This uses the stated sphere calculation, not a general
distributivity formula for a Hopf invariant of a composite.
The low tensor cells, equivalently the twice-suspended Ganea
calculation in D.2.49, give
\[
                 Z^{(9)}\simeq C_1^9\vee M(2,8)\vee S^9.
 \tag{D.2.60e}
\]
The map $M(2,7)\to C_1^9$ is injective on $\pi_{14}$.
Here an explicit unstable argument is necessary.  For
$f=i_7\eta_7:S^8\to M(2,7)$, the fiber of the pinch
$C_1^9\to S^9$ is the relative James construction.  Its
second stage attaches the cone on $M(2,7)\wedge S^7$ by
$[\mathrm{id},f]$; see
\cite[Lemmas~2.1 and 2.3--2.5, pp.~4--5]{FGHJinZhu2024}.
Its first new cell has dimension fifteen and attaches by
\[
 [i_7,i_7\eta_7]=i_7[\iota_7,\iota_7]\eta_{13}=0,
 \tag{D.2.60f}
\]
since $S^7$ is an $H$-space.  The next cell has dimension
sixteen; later word lengths start in dimension twenty-three.
Thus $M(2,7)$ maps isomorphically to the fiber on $\pi_{14}$.
The remaining boundary is zero: suspension
$\pi_{14}S^8\to\pi_{15}S^9$ is onto, and its image under the
fiber boundary is $i_7\eta_7\nu_8^2=0$.
Indeed $\eta_7\nu_8^2$ stabilizes to zero, while
$\pi_{14}S^7=\mathbb Z/8\{\sigma'\}$ stabilizes injectively
by $\sigma'\mapsto2\sigma$.  This verifies the hypotheses
of the cited relative-fiber exact sequence in the required
degree.  The exact Moore relation
$i_7\sigma'=2l_4\ne0$ has order two
\cite[Lemma~2.8(ii)]{EH5MorisugiMukai2000}.
Consequently the bottom class is nonzero of order two in
\textup{(D.2.60e)}.

\textbf{Literal integral cells and complete attaching matrices.}
Use the fixed mod-two labels
$a_2,b_3,c_4,d_5,e_5,f_6,x_6$ of $B$, where $d,e,f,x$
are dual to $w,uv,u^3,v^2$.  After the parity-preserving
integral normalization used in D.2.18 and D.2.22, choose chains
\[
 \partial b=2a,\quad \partial c=0,\quad
 \partial d=\partial e=4c,\quad
 \partial x=2(e-d),\quad \partial f=0.
 \tag{D.2.60g}
\]
The last cycle can also be chosen from the spherical class
of D.2.17.  Ordered tensor labels below suppress the shift
by three.  In the nine-stage set
\[
 B_8=ab,\quad e_8=ba-ab,\quad
 D_9=ac,\quad r_9=ca-ac,\quad m_9=bb-ac.
 \tag{D.2.60h}
\]
These give respectively the Chang, Moore and free factors
in \textup{(D.2.60e)}.  Write $\gamma,\delta\in\pi_9C_1^9$
for the free class with Hurewicz image $2D$ and the order-two
class detected by $\operatorname{Sq}^2$ on $B_8$.
The other possible torsion component is $e_8\eta_8$.

The literal ten-cycles are
$AD=ad-2bc$, $AE=ae-2bc$, $DA=da-2cb$, $EA=ea-2cb$;
the remaining cells are $BC=bc$ and $CB=cb$.
Their boundaries and second-square detectors determine all
components of their attaching maps:
\[
\begin{array}{c|rrrrrr}
\text{cell}&AD&AE&DA&EA&BC&CB\\ \hline
\text{attachment}&\delta&0&\delta+e_8\eta_8&0&
 \gamma+\delta+e_8\eta_8&\gamma+2r+\delta
\end{array}
\tag{D.2.60i}
\]
For example $\operatorname{Sq}_*^2(ad)=ab$ and
$\operatorname{Sq}_*^2(da)=ba$; the corresponding values on
$bc,cb$ are $ba,ab$.  There is no further component in
$\pi_9$ of the wedge.  The explicit cell changes
\[
 Q=AD,\quad V=DA-AD,\quad G=BC-DA,\quad
 H=CB-BC+DA-AD
\]
isolate $\delta,e_8\eta_8,\gamma,2r$.  Define
$\mathrm{Cb}_j=C_1^j\cup_\delta e^{j+1}\cup_\gamma e^{j+1}$.
We obtain an actual splitting
\[
 Z^{(10)}\simeq\mathrm{Cb}_9\vee C_1^{10}\vee M(2,9)
                         \vee S^{10}_{AE}\vee S^{10}_{EA}.
 \tag{D.2.60j}
\]
The complete degree-ten groups of these factors have free
generators $\kappa,\gamma',AE,EA$, with
$h(\kappa)=2Q$, $h(\gamma')=2V$, and three order-two
generators $v=i_7\nu_7,\delta',r\eta_9$.
Their independent torsion detectors are respectively
$\operatorname{Sq}^4$ on $aa$ and $\operatorname{Sq}^2$
on $m,r$.  Normalize $\kappa$ to have zero fourth-square
detector; its second-square detector is $D$.
These follow from the same Chang cofiber sequences and
Adem calculation as in D.2.38, shifted once.

Here is the entire signed eleven-cell matrix.  It records
the literal cellular boundary as well as the actual map;
thus no higher-exponent conclusion is being inferred from
the first Bockstein alone.
\[
\begin{array}{c|c|c}
\text{cell}&\partial&\text{attaching map}\\ \hline
AF&0&0\\
AX&2(AE-AD)&-\kappa+2AE\\
FA&0&0\\
XA&2(EA-DA)&-\kappa-\gamma'+2EA+r\eta_9\\
BD&2AD&\kappa+\delta'\\
BE&2AE&2AE\\
DB&-2DA&-\kappa-\gamma'+\delta'\\
EB&-2EA&-2EA\\
CC&0&v+r\eta_9
\end{array}
\tag{D.2.60k}
\]
In this order the second-square vectors are
$0,D,0,D+r,D+m,0,D+m,0,r$; only $CC$ has a nonzero
fourth square, namely $aa$.  Together with the integral
boundaries these are complete detectors in $\pi_{10}$.
First isolate $2AE,2EA$; then $DB+BD$ isolates $-\gamma'$,
and $AX+BD-BE$ isolates $\delta'$.  The remaining rows
isolate $\kappa,r\eta_9,v$ and two zero maps.  Consequently
\[
\begin{gathered}
 Z^{(11)}\simeq J\vee\mathrm{Cb}_{10}\vee C_1^{11}
                 \vee M(2,10)^{\vee2}\vee(S^{11})^{\vee2},\\
 J=\mathrm{Cb}_9\cup_\kappa e^{11}_K\cup_v e^{11}_N.
\end{gathered}
\tag{D.2.60l}
\]
These ordinary cell slides take place within the stable
range through dimension eleven and fix the bottom $S^7$.
In particular there is an actual retraction onto $J$.

\textbf{Every degree-eleven class of the bottom block.}
The Chang cofiber sequence gives
\[
\begin{gathered}
 \pi_{10}\mathrm{Cb}_9=\mathbb Z_{(2)}\{\kappa\}
                                  \oplus\mathbb Z/2\{v\},\\
 0\longrightarrow\mathbb Z/2\{\omega\}
 \longrightarrow\pi_{11}\mathrm{Cb}_9
 \longrightarrow\mathbb Z/2\{G\eta_{10}\}\longrightarrow0,
 \qquad\kappa\eta_{10}=\omega.
\end{gathered}
\tag{D.2.60m}
\]
For completeness, before the two ten-cells the relevant
Chang group is $\mathbb Z/4\{t\}$, with pinch coordinates
$(2\nu_8,\eta_9^2)$ and $\delta\eta_9^2=2t$.
The $\delta$ cell leaves its order-two image $\omega$;
the $\gamma$ cell supplies the displayed $G\eta_{10}$
quotient.  This is the cofiber calculation of D.2.38,
in degrees shifted once.  Attaching $\kappa$ kills
$\omega$ and supplies no free lift, because $\kappa$ has
infinite order.  Its remaining group is
$\mathbb Z/2\{\zeta\}$.  Finally $v$ has order two
and $v\eta_{10}=0$, so its eleven-cell gives
\[
 \pi_{11}J=\mathbb Z_{(2)}\{p\}\oplus\mathbb Z/2\{\zeta\},
 \qquad h(p)=2N.
 \tag{D.2.60n}
\]
The second-square detector of $\zeta$ on $G_{10}$ is one.

\textbf{A quotient detects the torsion composite.}
Collapsing the bottom gives the actual equivalence
\[
 \mathrm{Cb}_9/S^7\simeq C(\eta_8)\vee M(2,9).
 \tag{D.2.60o}
\]
Their degree-eleven groups are respectively
$\mathbb Z/4\{i_8\nu_8\}$ and
$\mathbb Z/4\{\widetilde\eta_9\}$.
The Moore coordinate of $\zeta$ is odd because its pinch
on $G$ is $\eta_{10}$.  The connecting map to $S^8$
sends the first generator to $2\nu_8$ and the second
to an odd multiple of $2\nu_8$, by the ordinary bracket
$\{\eta_8,2,\eta_9\}=\{2\nu_8,6\nu_8\}$.
Exactness forces its $C(\eta_8)$ coordinate to be odd too.
Therefore $\zeta\nu_{11}$ has the nonzero coordinate
$i_8\nu_8^2$ in that factor.  The fourth stable stem is
zero, so the $\eta$ cell does not kill this bottom
$\nu^2$; the Moore coordinate has zero composite by
the fifth stable stem of the mod-two Moore spectrum.
Adding the two eleven-cells to obtain $J/S^7$ still
preserves this coordinate: these quotient spaces are
seven-connected, and the possible kernel on $\pi_{14}$
comes from $\pi_{14}(S^{10}\vee S^{10})=0$.
In particular
\[
 \zeta\nu_{11}\text{ has nonzero image in }\pi_{14}(J/S^7).
 \tag{D.2.60p}
\]

\textbf{The free composite and the full twelve-cell kernel.}
There is an actual map $j:C(\nu_7)\to J$ on the bottom
and the $N$ cell.  The unique $v_8\in\pi_{11}C(\nu_7)$
of top degree eight exists by its stable-range cofiber
sequence: $\pi_{11}S^7=0$ and $\nu_7$ has order eight.
Both $jv_8$ and $4p$ have integral Hurewicz image $8N$.
Their difference could only be $\zeta$, but both have
zero second-square detector on $G_{10}$.  For $4p$ use
the degree-four cone map; for $jv_8$ use naturality and
$H^{10}C(\nu_7)=0$.  Hence $4p=jv_8$ exactly.
The ordinary coextension formula, with its full
indeterminacy, now gives
\[
 4p\nu_{11}=i_7\{\nu_7,8\iota_{10},\nu_{10}\}
            =i_7(4\sigma')=0.
 \tag{D.2.60q}
\]
The singleton bracket is the one proved in D.2.55 from
the ordinary $\sigma'''$ bracket and its two suspensions;
see also \cite[Section~2.5 and accompanying program]
{FGHStricklandToda}.  Its indeterminacy is
$\nu_7\pi_{14}S^{10}+\pi_{11}S^7\nu_{11}=0$.
The last equality in \textup{(D.2.60q)} is actual:
$2i_7=0$ and $[i_7,i_7]=0$ because $S^7$ is an
$H$-space, so the cyclic Hilton correction vanishes.

If $(mp+\epsilon\zeta)\nu_{11}$ lies in the old bottom
subgroup, pinching $N$ gives $2m\nu_{11}=0$ in
$\pi_{14}S^{11}=\mathbb Z/8\{\nu_{11}\}$.
Thus $4\mid m$, and \textup{(D.2.60q)} kills that part.
Equation \textup{(D.2.60p)} then forces $\epsilon=0$.
We have proved the complete intersection statement
\[
 \bigl(\pi_{11}J\circ\nu_{11}\bigr)
                    \cap\langle i_7\sigma'\rangle=0.
 \tag{D.2.60r}
\]
For the actual attachment of all twelve-cells, the
Blakers--Massey comparison from the wedge of attaching
$S^{11}$ spheres to the homotopy fiber is
sixteen-connected.  Thus its entire kernel on old
$\pi_{14}$ consists of the $\nu_{11}$ composites of
the attaching maps.  Retracting \textup{(D.2.60l)} to
$J$ and applying \textup{(D.2.60r)} excludes a bottom
relation.  The same comparison is fourteen-connected
for the ten-cells and fifteen-connected for the
eleven-cells; their possible contributions vanish by
$\pi_{14}S^9=\pi_5^s\mathbb S=0$ and
$\pi_{14}S^{10}=\pi_4^s\mathbb S=0$.
Together with the nine-stage argument, this proves
\textup{(D.2.60b)}.  No comparison in this proof
replaces the actual degree-fourteen homotopy group by
its stabilization, and no higher-cell conclusion follows.
\end{proof}
